\begin{quote}\itshape
What is the smallest timing rule that can turn many independent flash cycles into
one coherent display?
\end{quote}



\title{\textbf{\textcolor{titleblue}{Fireflies in Unison}}\\[0.3em]
       \large How Simple Timing Rules Create a Collective Rhythm}
\author{Yoshiki Kuramoto\textsuperscript{1} and Steven H. Strogatz\textsuperscript{2}\\[0.3em]
        {\small \textsuperscript{1}Kyoto, Japan \qquad
         \textsuperscript{2}Ithaca, New York, USA}}
\date{}

\begin{document}
\maketitle




The three parts of eq.~\eqref{eq:kuramoto} have direct meanings:
\begin{description}[style=nextline]
    \item[Internal clock] $\omega_i$ is the rate the firefly would follow alone.
    \item[Coupling strength] $K \ge 0$ measures how strongly flashes influence one
          another.
    \item[Timing correction] $\sin(\theta_j - \theta_i)$ tells firefly $i$ whether
          another phase is ahead or behind.
\end{description}




\begin{table}[ht]
    \centering
    \caption{How a neighboring firefly affects firefly $i$.}
    \label{tab:sine}
    \begin{tabular}{l l l >{\raggedright\arraybackslash}p{4.3cm}}
        \toprule
        \textbf{Relative position} & \textbf{Phase difference} &
        \textbf{Sine value} & \textbf{Effect on firefly $i$} \\
        \midrule
        \rowcolor{green!10}
        Neighbor ahead  & $0 < \theta_j-\theta_i < \pi$  &
            \textbf{\textcolor{green!55!black}{Positive}} & Speeds up its phase. \\
        \rowcolor{gray!15}
        Phases aligned  & $\theta_j-\theta_i = 0$        &
            \textbf{Zero} & Makes no correction. \\
        \rowcolor{red!10}
        Neighbor behind & $-\pi < \theta_j-\theta_i < 0$ &
            \textbf{\textcolor{red!70!black}{Negative}} & Slows down its phase. \\
        \rowcolor{gray!15}
        Opposite phases & $\theta_j-\theta_i = \pi$      &
            \textbf{Zero} & Produces no preferred direction of correction. \\
        \bottomrule
    \end{tabular}
\end{table}



itemize -- unordered
enumerate -- ordered 



\begin{algorithm}[ht]
    \caption{Simulate the Kuramoto model}
    \begin{algorithmic}[1]
        \Require Phases $\theta_i$, frequencies $\omega_i$, coupling $K$, step
                 size $h$, and number of steps $M$
        \For{$m \gets 1$ \textbf{to} $M$}
            \For{each firefly $i$}
                \State $q_i \gets \displaystyle\sum_{j=1}^{N}\sin(\theta_j-\theta_i),
                        \quad v_i \gets \omega_i + \frac{K}{N} q_i$
            \EndFor
            \State Update all phases simultaneously:
                   $\theta_i \gets (\theta_i + h v_i) \bmod 2\pi$
        \EndFor
        \State \Return $\theta_1,\dots,\theta_N$
    \end{algorithmic}
\end{algorithm}